\documentclass[11pt,a4paper]{report}

%============================================================
% Manuale tecnico del progetto
%   UAV Test Generator -- (1+1)-ES con diversita' DTW
%
% Compilabile con: pdflatex documentation.tex  (x2)
% oppure           xelatex documentation.tex   (x2)
%============================================================

\usepackage[utf8]{inputenc}
\usepackage[T1]{fontenc}
\usepackage{lmodern}
\usepackage[italian]{babel}

\usepackage[margin=2.4cm]{geometry}
\usepackage{microtype}
\usepackage{amsmath,amssymb}
\usepackage{xcolor}
\usepackage{graphicx}
\usepackage{booktabs}
\usepackage{tabularx}
\usepackage{longtable}
\usepackage{array}
\usepackage{enumitem}
\usepackage{listings}
\usepackage{caption}
\usepackage{titlesec}
\usepackage{fancyhdr}
\usepackage[hidelinks]{hyperref}

%------------------------------------------------------------
% Colori
%------------------------------------------------------------
\definecolor{PolimiBlue}{RGB}{0,83,155}
\definecolor{CodeBg}{RGB}{246,246,246}
\definecolor{PseudoBg}{RGB}{240,246,240}
\definecolor{AsciiBg}{RGB}{248,248,252}
\definecolor{KwBlue}{RGB}{20,60,120}
\definecolor{CmGreen}{RGB}{0,110,60}
\definecolor{StrRed}{RGB}{150,30,30}
\definecolor{RuleGray}{RGB}{200,200,200}

\hypersetup{
  colorlinks=true,
  linkcolor=PolimiBlue,
  urlcolor=PolimiBlue,
  citecolor=PolimiBlue,
  pdftitle={Manuale Tecnico - UAV Test Generator},
  pdfauthor={Documentazione generata come tutoraggio tecnico}
}

%------------------------------------------------------------
% Intestazioni
%------------------------------------------------------------
\pagestyle{fancy}
\fancyhf{}
\fancyhead[L]{\small UAV Test Generator -- Manuale Tecnico}
\fancyhead[R]{\small\thepage}
\renewcommand{\headrulewidth}{0.4pt}

\titleformat{\chapter}[display]
  {\normalfont\bfseries\color{PolimiBlue}}
  {\filleft\Large\chaptertitlename\ \thechapter}
  {1ex}{\titlerule\vspace{1ex}\filright\Huge}

\setcounter{tocdepth}{2}
\setcounter{secnumdepth}{3}
\setlist[itemize]{leftmargin=1.5em,itemsep=2pt,topsep=3pt}
\setlist[enumerate]{leftmargin=1.8em,itemsep=2pt,topsep=3pt}
\renewcommand{\arraystretch}{1.25}
\newcolumntype{Y}{>{\raggedright\arraybackslash}X}
\newcolumntype{L}[1]{>{\raggedright\arraybackslash}p{#1}}

%------------------------------------------------------------
% Codice inline: \code{...}  (underscore da scrivere come \_)
%------------------------------------------------------------
\newcommand{\code}[1]{\texttt{#1}}

%------------------------------------------------------------
% Stili listings
%------------------------------------------------------------
\lstdefinestyle{py}{
  language=Python,
  basicstyle=\ttfamily\footnotesize,
  backgroundcolor=\color{CodeBg},
  frame=single,
  rulecolor=\color{RuleGray},
  breaklines=true,
  columns=fullflexible,
  keepspaces=true,
  showstringspaces=false,
  numbers=left,
  numberstyle=\tiny\color{gray},
  keywordstyle=\color{KwBlue}\bfseries,
  commentstyle=\color{CmGreen}\itshape,
  stringstyle=\color{StrRed},
  tabsize=4
}

\lstdefinestyle{bash}{
  language=bash,
  basicstyle=\ttfamily\footnotesize,
  backgroundcolor=\color{CodeBg},
  frame=single,
  rulecolor=\color{RuleGray},
  breaklines=true,
  columns=fullflexible,
  showstringspaces=false,
  keywordstyle=\color{KwBlue}\bfseries,
  commentstyle=\color{CmGreen}\itshape,
  stringstyle=\color{StrRed}
}

% Pseudocodice (nessun linguaggio, parole chiave evidenziate a mano)
\lstdefinestyle{pseudo}{
  basicstyle=\ttfamily\footnotesize,
  backgroundcolor=\color{PseudoBg},
  frame=single,
  rulecolor=\color{RuleGray},
  breaklines=true,
  columns=fullflexible,
  keepspaces=true,
  mathescape=false,
  morekeywords={se,allora,altrimenti,per,ogni,finche,mentre,ritorna,funzione,fine,in,and,or,not,da,a,fino,restituisci,break,continua},
  keywordstyle=\color{PolimiBlue}\bfseries,
  commentstyle=\color{CmGreen}\itshape,
  tabsize=2
}

% Diagrammi ASCII (solo caratteri ASCII, nessun box-drawing Unicode)
\lstdefinestyle{ascii}{
  basicstyle=\ttfamily\scriptsize,
  backgroundcolor=\color{AsciiBg},
  frame=single,
  rulecolor=\color{RuleGray},
  columns=fullflexible,
  keepspaces=true,
  breaklines=false
}

%============================================================
\begin{document}
%============================================================

%------------------------------------------------------------
% Frontespizio
%------------------------------------------------------------
\begin{titlepage}
\centering
\vspace*{2cm}
{\Large Politecnico di Milano\par}
\vspace{0.3cm}
{\large Software Engineering for Automation\par}
\vfill
{\Huge\bfseries\color{PolimiBlue} Manuale Tecnico\par}
\vspace{0.6cm}
{\Large UAV Test Generator\par}
{\large Generazione di test basata su ricerca con\\
Strategia Evolutiva $(1+1)$ e diversità di traiettoria (DTW)\par}
\vfill
{\large Guida completa all'architettura, al codice e al flusso di esecuzione\par}
\vspace{1cm}
{\normalsize Documento di studio e comprensione del progetto\par}
\vfill
\end{titlepage}

\tableofcontents
\clearpage

%============================================================
\chapter{Introduzione}
%============================================================

\section{Scopo di questo documento}
Questo manuale è pensato come una guida di studio completa al progetto
\emph{UAV Test Generator}. L'obiettivo non è solo elencare cosa fa ogni
riga di codice, ma costruire una comprensione profonda:
\begin{itemize}
  \item \textbf{del problema} che il software risolve;
  \item \textbf{dell'architettura} e del perché è stata scelta;
  \item \textbf{del flusso di esecuzione} end-to-end (chi chiama cosa e quando);
  \item \textbf{degli algoritmi} che stanno al cuore del sistema, spiegati
        con pseudocodice passo per passo;
  \item \textbf{dei concetti teorici} necessari (strategie evolutive,
        Dynamic Time Warping, regola $1/5$ di Rechenberg) presentati
        senza darli per scontati.
\end{itemize}

Alla fine della lettura si deve essere in grado di comprendere,
modificare ed estendere il progetto senza dover ripartire dal codice
sorgente.

\section{Scopo del progetto}
Il software partecipa alla \emph{UAV Testing Competition}. Il compito è
\textbf{generare automaticamente scenari di test} (configurazioni di
ostacoli) che mettano in difficoltà il sistema di
\emph{obstacle avoidance} di un drone autonomo, portandolo a violare la
distanza di sicurezza o addirittura a collidere.

\begin{itemize}
  \item Il \textbf{sistema sotto test} (System Under Test, SUT) è
        \emph{PX4-Avoidance}: il modulo di evitamento ostacoli del
        firmware di volo PX4. Non viene modificato: è una scatola nera
        che riceve una missione e degli ostacoli, vola in simulazione e
        restituisce la traiettoria percorsa.
  \item Il \textbf{banco di prova} è \emph{Aerialist}, un framework che
        configura ed esegue simulazioni PX4/Gazebo e restituisce i log
        di volo.
  \item Il nostro software è il \textbf{generatore di test}: propone
        configurazioni di ostacoli, le fa simulare, osserva il risultato
        e usa questa informazione per proporre configurazioni sempre più
        pericolose.
\end{itemize}

Il criterio di qualità di un test è la \textbf{distanza minima}
$d_{\min}$ tra la traiettoria del drone e gli ostacoli durante il volo:
più è piccola, più il test è severo. Un buon insieme di test finale è
inoltre \emph{diversificato} (traiettorie di fallimento diverse tra
loro) e \emph{semplice} (pochi ostacoli, voli brevi).

\section{Idea centrale in una frase}
\begin{quote}
Il generatore usa una \textbf{Strategia Evolutiva $(1+1)$} che
minimizza la distanza minima drone--ostacolo, campionando gli ostacoli
lungo la rotta nominale del drone, e conserva in un \textbf{archivio}
solo i fallimenti con traiettorie sufficientemente diverse tra loro,
misurate con il \textbf{Dynamic Time Warping}.
\end{quote}

%============================================================
\chapter{Concetti teorici preliminari}
%============================================================

Prima di leggere il codice conviene fissare quattro concetti. Il resto
del manuale vi farà riferimento continuamente.

\section{Strategia Evolutiva \texorpdfstring{$(1+1)$}{(1+1)}}
Una \emph{Strategia Evolutiva} (Evolution Strategy, ES) è un algoritmo
di ottimizzazione ispirato all'evoluzione naturale. La variante
$(1+1)$ è la più semplice possibile:
\begin{itemize}
  \item si mantiene \textbf{una sola soluzione}, chiamata
        \emph{genitore} (\emph{parent});
  \item a ogni iterazione (\emph{generazione}) si crea \textbf{un solo
        figlio} (\emph{child}) applicando una \emph{mutazione} casuale
        al genitore;
  \item il figlio \textbf{sostituisce} il genitore solo se è migliore
        (selezione elitaria).
\end{itemize}
La notazione $(1+1)$ significa: 1 genitore $+$ 1 figlio, dai quali si
seleziona il migliore. È adatta quando ogni valutazione (qui: una
simulazione) è costosa, perché usa pochissime valutazioni per
generazione e occupa poca memoria.

\section{Fitness e minimizzazione}
La \emph{fitness} è la funzione obiettivo. Qui:
\[
f(x) = d_{\min}(x), \qquad f:\mathcal{X}\to\mathbb{R}_{\ge 0},
\]
dove $x$ è una configurazione di ostacoli e $d_{\min}(x)$ è la distanza
minima osservata durante il volo. Il problema è di
\textbf{minimizzazione}: valori più \emph{bassi} sono \emph{migliori}
(più pericolosi). Il numero di ostacoli \emph{non} entra nella fitness:
serve solo dopo, nel ranking di consegna.

\section{Regola \texorpdfstring{$1/5$}{1/5} di Rechenberg}
Quanto devono essere ampie le mutazioni? Se sono troppo piccole la
ricerca è lenta; se troppo grandi diventa casuale. La regola
\textbf{$1/5$ di Rechenberg} adatta automaticamente l'ampiezza del passo
$\sigma$ osservando il \emph{tasso di successo} (quante mutazioni recenti
hanno migliorato il genitore):
\begin{itemize}
  \item se più di $1/5$ ($20\%$) delle mutazioni recenti migliora
        $\Rightarrow$ la ricerca sta trovando terreno facile
        $\Rightarrow$ \textbf{aumenta} $\sigma$ (esplora più lontano);
  \item altrimenti $\Rightarrow$ sta faticando $\Rightarrow$
        \textbf{riduce} $\sigma$ (affina localmente).
\end{itemize}
È il classico compromesso \emph{esplorazione} vs \emph{sfruttamento}.

\section{Dynamic Time Warping (DTW)}
Due configurazioni di ostacoli diverse possono produrre la
\emph{stessa} manovra di evitamento. Ciò che conta come ``diverso'', ai
fini della competizione, è la \textbf{traiettoria}, non l'ostacolo. Il
\textbf{Dynamic Time Warping} è una distanza tra due serie temporali
(qui: due traiettorie $(x,y)$) che le allinea in modo elastico,
tollerando sfasamenti nel tempo. Se il DTW medio tra due traiettorie è
sotto una soglia, le si considera la \emph{stessa classe di
comportamento}.

Formalmente, date $A=(a_1,\dots,a_n)$ e $B=(b_1,\dots,b_m)$, si riempie
una matrice $D$ con la ricorrenza:
\[
D_{i,j} = \lVert a_i - b_j\rVert_2 +
\min\{D_{i-1,j},\; D_{i,j-1},\; D_{i-1,j-1}\},
\]
con $D_{0,0}=0$ e bordi a $+\infty$. Il valore finale $D_{n,m}$ viene
normalizzato sulla lunghezza dell'allineamento per ottenere una
deviazione media in metri.

%============================================================
\chapter{Architettura del progetto}
%============================================================

\section{Visione d'insieme}
Il sistema segue un'architettura a \textbf{pipeline modulare} con una
netta separazione tra:
\begin{enumerate}
  \item \textbf{logica di ricerca deterministica} (mutazione, selezione,
        scoring, DTW) --- pura, testabile, veloce;
  \item \textbf{interazione con il simulatore} --- costosa, lenta,
        potenzialmente instabile, isolata dietro un adattatore.
\end{enumerate}
Questa separazione (nota come \emph{Functional Core, Imperative Shell})
è la decisione architetturale più importante: permette di testare tutta
la logica senza avviare Gazebo, e confina il codice ``fragile'' in un
punto solo.

\section{Diagramma dei componenti}
\begin{lstlisting}[style=ascii]
                    +----------------------------+
   Operatore  --->  |  start.sh  (launcher)      |  avvio + auto-resume
                    +-------------+--------------+
                                  |  docker run
                                  v
                    +----------------------------+
                    |  cli.py  (entry point)     |  parsing + logging
                    +-------------+--------------+
                                  |  seleziona il generatore (env GENERATOR)
                                  v
      +---------------------------+----------------------------+
      |            ESGenerator  (es_generator.py)              |
      |  orchestrazione: valutazione, checkpoint, ciclo (1+1), |
      |  archivio, consegna                                    |
      +--+-------------+-------------+------------+-------------+
         |             |             |            |
         v             v             v            v
   +-----------+ +-----------+ +-----------+ +-----------+
   | config.py | |geometry.py| | models.py | |testcase.py|
   | costanti  | | funzioni  | |Individual | | adattatore|
   | parametri | | pure      | | + serial. | | Aerialist |
   +-----------+ +-----------+ +-----------+ +-----+-----+
                                                   |
                       (RandomGenerator eredita)   v
                                             +-----------+
                                             | Aerialist |
                                             | PX4/Gazebo|
                                             +-----+-----+
                                                   |
                                                   v
                                             +-----------+
                                             |Filesystem |
                                             | checkpoint|
                                             | all_fails |
                                             | consegna  |
                                             +-----------+
\end{lstlisting}

\section{Le due directory di output: \texttt{all\_fails} e \texttt{consegna}}
Concetto chiave per capire tutto il resto:
\begin{description}
  \item[\code{all\_fails/}] \emph{storia completa}. Ogni fallimento
    valido trovato viene salvato qui immediatamente e non viene mai
    cancellato. Serve come prova permanente e per ricostruire la
    consegna.
  \item[\code{consegna/}] \emph{sottoinsieme da valutare}. Solo i
    fallimenti selezionati (diversi tra loro per traiettoria, ordinati
    per qualità) vengono \emph{copiati} qui con nomi \code{rankNN\_*}.
\end{description}
La consegna è sempre ricostruibile da \code{all\_fails}; per questo si
\emph{copia} e non si \emph{sposta}.

%============================================================
\chapter{Struttura del repository}
%============================================================

\begin{longtable}{L{4.2cm} L{10cm}}
\toprule
\textbf{File / cartella} & \textbf{Ruolo}\\
\midrule
\endfirsthead
\toprule
\textbf{File / cartella} & \textbf{Ruolo}\\
\midrule
\endhead
\code{cli.py} & Punto di ingresso. Parsing argomenti, selezione del generatore, logging, gestione errori.\\
\code{config.py} & Costanti ufficiali e parametri configurabili via variabili d'ambiente.\\
\code{geometry.py} & Funzioni pure: geometria ostacoli, campionamento/mutazione geni, rotta, DTW, formule di punteggio.\\
\code{models.py} & Classe \code{Individual} (genoma + metriche) e sua serializzazione JSON.\\
\code{es\_generator.py} & Classe \code{ESGenerator}: valutazione, checkpoint, ciclo $(1+1)$-ES, archivio, consegna.\\
\code{random\_generator.py} & Baseline di ricerca casuale (sottoclasse di \code{ESGenerator}).\\
\code{testcase.py} & Adattatore verso Aerialist (esecuzione di una simulazione).\\
\code{plot\_official.py} & Generazione dei grafici ufficiali per i test salvati.\\
\code{start.sh} & Launcher Docker con auto-resume su crash.\\
\code{Dockerfile} & Immagine self-contained (estende l'immagine Aerialist).\\
\code{requirements.txt} & Dipendenze Python aggiuntive (numpy, matplotlib, pyyaml).\\
\code{case\_studies/} & Missioni di input (\code{.yaml}, \code{.plan}, \code{.ulg}) e immagini.\\
\code{generated\_tests/} & Output delle campagne: checkpoint, log, \code{all\_fails/}, \code{consegna/}.\\
\code{docs/uml/} & Diagrammi Mermaid: classi, flusso di esecuzione, sequenza.\\
\bottomrule
\end{longtable}

\section{Ordine di studio consigliato}
Per comprendere il progetto senza perdersi, consiglio questo ordine
(dal più semplice e ``foglia'' al più complesso e ``radice''):
\begin{enumerate}
  \item \code{config.py} --- il vocabolario: tutte le costanti e i
        parametri. Si legge in 10 minuti e serve per tutto il resto.
  \item \code{models.py} --- la struttura dati centrale
        (\code{Individual}). Piccolo e fondamentale.
  \item \code{geometry.py} --- le funzioni pure. Qui vivono gli algoritmi
        matematici (mutazione, DTW, trap-gene). Denso ma autonomo.
  \item \code{testcase.py} --- l'unico ponte verso il simulatore.
        Brevissimo.
  \item \code{es\_generator.py} --- il cuore. Mette insieme tutto:
        ciclo evolutivo, archivio, checkpoint, consegna.
  \item \code{random\_generator.py} --- si capisce in 5 minuti dopo il
        precedente (cambia solo due metodi).
  \item \code{cli.py} --- il punto d'ingresso, ora che si sa cosa
        orchestra.
  \item \code{plot\_official.py}, \code{start.sh}, \code{Dockerfile}
        --- contorno operativo.
\end{enumerate}

%============================================================
\chapter{Descrizione dei moduli e loro comunicazione}
%============================================================

\section{Grafo delle dipendenze}
\begin{lstlisting}[style=ascii]
        cli.py
        /     \
       v       v
 es_generator.py <---- random_generator.py  (eredita)
   |    |    |   \
   v    v    v    v
config geom model testcase ----> Aerialist ----> PX4/Gazebo
  ^     |    |
  |     v    v
  +--- config (geometry e models importano le costanti)

 plot_official.py ---> numpy / matplotlib / pyyaml   (usato da es_generator
                                                       solo via import "lazy")
\end{lstlisting}
Le frecce indicano ``dipende da / importa''. Da notare:
\begin{itemize}
  \item le dipendenze puntano \emph{verso l'interno}, verso moduli
        semplici (\code{config}, \code{geometry}, \code{models});
  \item \code{geometry.py} e \code{models.py} non conoscono
        \code{ESGenerator}: sono riutilizzabili e testabili in
        isolamento;
  \item solo \code{testcase.py} (e di riflesso \code{geometry} per la
        conversione ostacoli) tocca la libreria Aerialist.
\end{itemize}

\section{Tabella riassuntiva delle responsabilità}
\begin{longtable}{L{3.4cm} L{4.2cm} L{6.4cm}}
\toprule
\textbf{Modulo} & \textbf{Responsabilità} & \textbf{Cosa NON fa}\\
\midrule
\endfirsthead
\toprule
\textbf{Modulo} & \textbf{Responsabilità} & \textbf{Cosa NON fa}\\
\midrule
\endhead
\code{config.py} & Definire costanti e leggere le env. & Nessuna logica.\\
\code{geometry.py} & Calcoli geometrici, campionamento, mutazione, DTW, scoring. & Non gestisce stato, non simula, non fa I/O (tranne leggere la rotta).\\
\code{models.py} & Rappresentare un candidato e serializzarlo. & Non calcola nulla.\\
\code{testcase.py} & Eseguire una simulazione via Aerialist. & Non conosce la ricerca.\\
\code{es\_generator.py} & Orchestrare l'intera campagna. & Non implementa la fisica né i calcoli puri (li delega).\\
\code{cli.py} & Avviare e collegare i pezzi. & Nessun algoritmo.\\
\bottomrule
\end{longtable}

%============================================================
\chapter{Analisi dettagliata di ogni file}
%============================================================

%------------------------------------------------------------
\section{\texttt{config.py} --- costanti e parametri}
\label{sec:config}
%------------------------------------------------------------

\subsection{Ruolo}
È il ``vocabolario'' del progetto: tutti i numeri magici stanno qui,
ciascuno leggibile da variabile d'ambiente con un valore di default.
Questo è cruciale perché \textbf{il valutatore lancia il generatore
senza impostare alcuna variabile}: i default \emph{sono} l'algoritmo che
gira in valutazione.

\subsection{Elementi importanti}
\begin{itemize}
  \item \textbf{Range ufficiali} (spazio legale degli ostacoli):
  \[
  2\le l\le 20,\quad 2\le w\le 20,\quad 10\le h\le 25,
  \]
  \[
  -40\le x\le 30,\quad 10\le y\le 40,\quad 0^\circ\le r\le 90^\circ.
  \]
  Codificati in \code{MIN\_SIZE/MAX\_SIZE} e \code{MIN\_POS/MAX\_POS}, da
  cui derivano i vettori \code{GENE\_LOW} e \code{GENE\_HIGH}.
  \item \textbf{Gene}: un ostacolo è il vettore
  $[l,w,h,x,y,r]$ (la quota $z$ è sempre $0$).
  \item \textbf{Costanti di esito archivio}: \code{ARCH\_NONE},
  \code{ARCH\_NEW}, \code{ARCH\_IMPROVED}, \code{ARCH\_REDUNDANT}
  (valori $0,1,2,3$), usate come segnale di novità per la ricerca.
  \item \textbf{Seed}: se la variabile \code{SEED} è impostata, si
  inizializza il PRNG di Python per rendere ripetibile il campionamento
  (le simulazioni restano comunque stocastiche).
\end{itemize}

\subsection{Catalogo completo dei parametri}
\begin{longtable}{L{3.6cm} L{1.8cm} L{8.6cm}}
\toprule
\textbf{Parametro} & \textbf{Default} & \textbf{Significato}\\
\midrule
\endfirsthead
\toprule
\textbf{Parametro} & \textbf{Default} & \textbf{Significato}\\
\midrule
\endhead
\code{GENERATOR} & \code{es} & Sceglie \code{es} (evolutivo) o \code{random} (baseline).\\
\code{N\_RUNS} & 1 & Simulazioni per valutazione durante la ricerca.\\
\code{TEST\_TIMEOUT} & 500 & Secondi massimi per singola simulazione.\\
\code{SIM\_RETRIES} & 3 & Tentativi entro una chiamata \code{\_run\_once}.\\
\code{CONFIRM\_BELOW} & 1.0 & Soglia (m) sotto cui scatta la conferma anti-rumore.\\
\code{CONFIRM\_EXTRA\_RUNS} & 1 & Simulazioni extra di conferma.\\
\code{ADAPTIVE} & 1 & Campionamento lungo la rotta ($0$ = box fisso).\\
\code{ROUTE\_JITTER} & 3.0 & Deviazione std (m) del jitter attorno alla rotta.\\
\code{DEST\_TOLERANCE} & 5.0 & Raggio (m) per considerare la destinazione raggiunta.\\
\code{COMPLETE\_MIN\_FRAC} & 0.5 & Frazione di run ``arrivati'' per dire ``completato''.\\
\code{SIGMA\_INIT} & 0.15 & Passo di mutazione iniziale (frazione del range).\\
\code{SIGMA\_MIN} & 0.02 & Limite inferiore del passo.\\
\code{SIGMA\_MAX} & 0.4 & Limite superiore del passo.\\
\code{SOFT\_FAIL} & 1.5 & Distanza (m) sotto cui è un fallimento.\\
\code{STAGNATION\_LIMIT} & 15 & Generazioni senza miglioramento prima del restart.\\
\code{INIT\_SEEDS} & 5 & Candidati iniziali valutati per lo start.\\
\code{RESTART\_CANDIDATES} & 12 & Candidati generati a ogni restart (si tiene il meno simile).\\
\code{INTENSIFY\_PROB} & 0.4 & Probabilità di intensificare quando le zone sono coperte.\\
\code{ARCHIVE\_CAP} & 20 & Capienza dell'archivio di consegna.\\
\code{DUP\_IOU} & 0.95 & Soglia IoU per duplicati (fallback senza traiettoria).\\
\code{Y\_BUCKET\_SIZE} & 10 & Larghezza del bucket lungo $y$ per la diversità spaziale.\\
\code{DTW\_THRESHOLD} & 1.5 & Deviazione DTW (m) sotto cui due traiettorie sono ``uguali''.\\
\code{TRAJ\_N} & 100 & Punti per traiettoria ricampionata (DTW è $O(n^2)$).\\
\code{POS\_STEP\_FRAC} & 0.6 & Frazione dell'estensione rotta per il passo posizionale.\\
\code{POS\_STEP\_MIN} & 3.0 & Passo posizionale minimo (m).\\
\code{POS\_STEP\_MAX} & 8.0 & Passo posizionale massimo (m).\\
\code{TRAP\_PROB} & 0.7 & Prob. che l'ostacolo aggiunto sia mirato (trap-gene) vs casuale.\\
\code{TRAP\_JITTER} & 1.0 & Dispersione (m) del trap-gene attorno al varco.\\
\code{TRAP\_SPAN\_MULT} & 3.0 & Lato lungo del trap $=$ mult.\ $\times$ varco misurato.\\
\code{TRAP\_SIZE\_FRAC} & 0.35 & Tetto sul lato lungo, come frazione dello span ufficiale.\\
\bottomrule
\end{longtable}

\begin{quote}
\textbf{Nota.} Il passo di mutazione posizionale è \emph{disaccoppiato}
dai limiti legali: viene tarato sul corridoio della rotta (vedi
\code{\_set\_step\_scale}). Senza questo, ogni figlio ``salterebbe via''
dalla rotta e la ES degenererebbe in ricerca casuale.
\end{quote}

%------------------------------------------------------------
\section{\texttt{models.py} --- la classe \texttt{Individual}}
\label{sec:models}
%------------------------------------------------------------

\subsection{Ruolo}
Definisce l'oggetto centrale del dominio: un \emph{candidato}. È
volutamente ``povero'': contiene solo valori semplici (numeri, liste) e
un riferimento temporaneo al \code{TestCase}. Questo lo rende
serializzabile in JSON e disaccoppia lo stato di ricerca dagli oggetti
pesanti di Aerialist.

\subsection{Attributi}
\begin{longtable}{L{3.4cm} L{3cm} L{7.6cm}}
\toprule
\textbf{Attributo} & \textbf{Tipo} & \textbf{Scopo}\\
\midrule
\endfirsthead
\toprule
\textbf{Attributo} & \textbf{Tipo} & \textbf{Scopo}\\
\midrule
\endhead
\code{genes} & lista di liste & Da uno a tre geni $[l,w,h,x,y,r]$.\\
\code{fitness} & float / None & Obiettivo di ricerca (distanza minima).\\
\code{min\_distance} & float & Distanza minima osservata.\\
\code{completed} & bool & Il drone ha raggiunto la destinazione?\\
\code{test\_case} & TestCase / None & Miglior artefatto d'esecuzione (non serializzato).\\
\code{fail\_base} & str / None & Percorso base dei file del fallimento su disco.\\
\code{avg\_point} & float / None & Punto medio (Formula 1) sui run.\\
\code{avg\_flight\_min} & float / None & Durata media di volo (minuti), per il ranking.\\
\code{traj} & lista / None & Traiettoria media ricampionata, per il DTW.\\
\bottomrule
\end{longtable}

\subsection{Metodi e serializzazione}
\begin{itemize}
  \item \code{obstacles()} converte i geni in oggetti \code{Obstacle} di
        Aerialist (chiamando \code{\_gene\_to\_obstacle} di
        \code{geometry.py}).
  \item \code{\_ind\_to\_dict} / \code{\_ind\_from\_dict} definiscono un
        \emph{confine di serializzazione} esplicito: escludono gli
        oggetti Aerialist, rappresentano l'infinito con \code{None} e
        supportano fallback retro-compatibili per checkpoint vecchi.
\end{itemize}
Questo è di fatto il pattern \textbf{Memento}: lo stato logico viene
catturato in una forma semplice e ripristinabile.

%------------------------------------------------------------
\section{\texttt{geometry.py} --- funzioni pure}
\label{sec:geometry}
%------------------------------------------------------------

\subsection{Ruolo}
Contiene tutta la matematica del progetto sotto forma di funzioni
(quasi) pure: nessuno stato di classe, facilmente testabili. Tre valori
a livello di modulo sono l'eccezione ``mutabile'': \code{\_ROUTE\_XY}
(punti della rotta), \code{\_DESTINATION} (destinazione), e
\code{GENE\_STEP} (passi di mutazione tarati sulla rotta).

\subsection{Gruppi di funzioni}
\begin{description}
  \item[Conversione e validità geometrica]
    \code{\_gene\_to\_obstacle}, \code{\_obstacle\_corners},
    \code{\_boxes\_overlap} (Teorema degli Assi Separatori),
    \code{\_overlaps\_any}, \code{\_fit\_center} (rientra il box ruotato
    nei limiti legali), \code{\_repair\_overlaps}.
  \item[Rotta e completamento]
    \code{\_load\_route} (legge la rotta dal \code{.ulg}),
    \code{\_reached\_destination}, \code{\_set\_step\_scale}.
  \item[Campionamento e mutazione]
    \code{\_random\_gene}, \code{\_mutate\_gene},
    \code{\_non\_overlapping\_gene}, \code{\_trap\_gene}.
  \item[Scoring ufficiale]
    \code{\_point} (Formula 1), \code{\_ranking\_score} (Formula 2).
  \item[Traiettorie e DTW]
    \code{\_resample\_xy}, \code{\_dtw\_mean}, \code{\_traj\_gap}.
\end{description}

\subsection{Il Teorema degli Assi Separatori (SAT)}
Per sapere se due box ruotati si sovrappongono, \code{\_boxes\_overlap}
usa il SAT: due poligoni convessi \emph{non} si sovrappongono se e solo
se esiste un asse su cui le loro proiezioni non si intersecano. Si
provano come assi le normali ai lati di entrambi i box; se su almeno un
asse le proiezioni sono disgiunte, i box sono separati.

\subsection{La funzione \texttt{\_fit\_center}}
Un box ruotato di $r$ gradi ha un \emph{bounding box} allineato agli
assi più grande del box stesso. \code{\_fit\_center} calcola le mezze
estensioni di questo bounding box e sposta il centro $(x,y)$ quel tanto
che basta perché \emph{tutto} il box (non solo il centro) resti nell'area
legale:
\[
\text{half}_x = \Big|\tfrac{l}{2}\cos r\Big| + \Big|\tfrac{w}{2}\sin r\Big|,
\qquad
\text{half}_y = \Big|\tfrac{l}{2}\sin r\Big| + \Big|\tfrac{w}{2}\cos r\Big|.
\]

\subsection{Le due formule ufficiali}
\textbf{Formula 1} (\code{\_point}): punteggio discreto di una singola
simulazione dalla distanza minima.
\[
p(d)=
\begin{cases}
5, & d<0.25,\\
2, & 0.25\le d<1.0,\\
1, & 1.0\le d<1.5,\\
0, & d\ge 1.5.
\end{cases}
\]
\textbf{Formula 2} (\code{\_ranking\_score}): punteggio di consegna,
usato \emph{solo} per ordinare/selezionare, mai nella fitness di ricerca.
\[
S(x)=\frac{10\,\overline{p}(x)}{n(x)^2\,\overline{t}(x)},
\]
dove $\overline{p}$ è il punto medio, $n$ il numero di ostacoli,
$\overline{t}$ il tempo medio di volo. Premia fallimenti severi, pochi
ostacoli (penalità quadratica) e voli brevi.

\subsection{DTW in pseudocodice}
La funzione \code{\_dtw\_mean} implementa la ricorrenza e poi normalizza
sulla lunghezza dell'allineamento tramite \emph{backtracking}.
\begin{lstlisting}[style=pseudo]
funzione dtw_mean(A, B):
    n, m = lunghezza(A), lunghezza(B)
    se n == 0 or m == 0: ritorna +infinito

    # matrice dei costi euclidei tra ogni coppia di punti
    cost[i][j] = || A[i] - B[j] ||

    # riempimento programmazione dinamica
    acc = matrice (n+1) x (m+1) inizializzata a +infinito
    acc[0][0] = 0
    per i da 1 a n:
        per j da 1 a m:
            acc[i][j] = cost[i-1][j-1]
                        + min(acc[i-1][j], acc[i][j-1], acc[i-1][j-1])

    # backtracking per contare i passi dell'allineamento
    i, j, passi = n, m, 0
    mentre i > 0 and j > 0:
        passi = passi + 1
        vai alla cella vicina con acc minimo (diagonale/su/sinistra)
    passi = passi + i + j
    ritorna acc[n][m] / max(passi, 1)   # deviazione media in metri
\end{lstlisting}
Il costo è $O(n\cdot m)$: per questo le traiettorie sono ricampionate a
\code{TRAJ\_N}$=100$ punti fissi (\code{\_resample\_xy}), mantenendo il
calcolo trattabile.

\subsection{Il trap-gene: l'operatore ``trappola''}
\code{\_trap\_gene} è l'operatore più originale del progetto. Quando la
mutazione decide di \emph{aggiungere} un ostacolo, invece di piazzarlo a
caso lo mette \emph{nel varco che il drone ha usato per scappare}.
\begin{lstlisting}[style=pseudo]
funzione trap_gene(individuo, ostacoli_esistenti):
    se l'individuo non ha traiettoria: ritorna None

    # 1. trova il punto della traiettoria piu' vicino agli ostacoli:
    #    e' il "varco" attraverso cui il drone e' passato
    per ogni punto p della traiettoria:
        d = distanza minima di p dagli ostacoli esistenti
    (tx, ty) = punto con d minima
    best_d = quella distanza minima (ampiezza del varco)

    # 2. stima la direzione di moto locale (tangente discreta)
    heading = angolo tra il punto successivo e quello precedente
    perp = heading + 90 gradi   # orientamento perpendicolare al moto

    # 3. dimensiona il lato lungo sul varco misurato (con un tetto)
    lato_lungo = clamp(TRAP_SPAN_MULT * best_d, minimo, cap)

    # 4. prova a piazzare il box perpendicolare, con jitter, senza overlap
    per alcuni tentativi:
        costruisci gene [l, w, h, x=tx+rumore, y=ty+rumore, r=perp]
        se non si sovrappone agli esistenti: ritorna il gene
    ritorna None
\end{lstlisting}
L'orientamento \emph{perpendicolare} al moto massimizza quanto il nuovo
ostacolo ostruisce il varco, invece di affiancarlo inutilmente.

%------------------------------------------------------------
\section{\texttt{testcase.py} --- adattatore verso Aerialist}
\label{sec:testcase}
%------------------------------------------------------------

\subsection{Ruolo}
È l'unico ponte diretto verso il simulatore, e implementa il pattern
\textbf{Adapter}: nasconde la scelta dell'agente (locale, Docker, K8s) e
offre un'interfaccia minimale.

\begin{lstlisting}[style=py]
class TestCase:
    def __init__(self, casestudy, obstacles):
        self.test = copy.deepcopy(casestudy)      # missione immutabile
        self.test.simulation.obstacles = obstacles # inietta gli ostacoli

    def execute(self):                # avvia la simulazione
        agent = <Local|Docker|K8s>Agent(self.test)
        self.test_results = agent.run()
        self.trajectory = self.test_results[0].record
        self.log_file   = self.test_results[0].log_file
        return self.trajectory

    def get_distances(self):          # distanza traiettoria-ostacolo
        return [self.trajectory.min_distance_to_obstacles([o])
                for o in self.test.simulation.obstacles]

    def save_yaml(self, path):
        self.test.to_yaml(path)
\end{lstlisting}

\subsection{Punti importanti}
\begin{itemize}
  \item Il \code{deepcopy} garantisce che la missione originale non venga
        mai modificata: si sostituisce \emph{solo} la lista di ostacoli.
  \item \code{get\_distances} restituisce una distanza per ostacolo; il
        chiamante ne prende il minimo (la fitness).
  \item L'agente è scelto dalla variabile \code{AGENT} (default
        \code{local} nel deployment).
\end{itemize}

%------------------------------------------------------------
\section{\texttt{es\_generator.py} --- il cuore del sistema}
\label{sec:esgen}
%------------------------------------------------------------

\subsection{Ruolo}
\code{ESGenerator} è al tempo stesso il \emph{servizio applicativo} e la
\emph{macchina a stati} del sistema. Coordina sei aree di
responsabilità:
\begin{enumerate}
  \item \textbf{valutazione}: retry, timeout, budget, run ripetuti,
        aggregazione;
  \item \textbf{bootstrap rotta}: simulazione senza ostacoli se manca la
        rotta;
  \item \textbf{checkpoint}: salvataggio/ripristino atomico dello stato;
  \item \textbf{ricerca}: seeding, selezione $(1+1)$, adattamento del
        passo, stagnazione, restart, intensificazione;
  \item \textbf{archivio}: dump dei fallimenti, deduplicazione per
        traiettoria, insieme limitato di rappresentanti;
  \item \textbf{consegna}: ordinamento per punteggio, separazione DTW
        finale, copia idempotente.
\end{enumerate}

\subsection{Valutazione: \texttt{\_run\_once} e \texttt{\_evaluate}}
\code{\_run\_once} esegue \emph{una} simulazione con retry e timeout,
consumando budget a ogni tentativo:
\begin{lstlisting}[style=pseudo]
funzione _run_once(ostacoli):
    per tentativo da 1 a SIM_RETRIES:
        se budget_used >= total_budget: ritorna None
        tc = TestCase(case_study, ostacoli)
        budget_used += 1                     # il budget si consuma SEMPRE
        imposta allarme di TEST_TIMEOUT secondi
        prova:
            tc.execute()
            d = min(tc.get_distances())
            annulla allarme; ritorna tc con min_distance = d
        se timeout o errore:
            annulla allarme; logga; passa al tentativo successivo
    ritorna None
\end{lstlisting}
\code{\_evaluate} chiama \code{\_run\_once} per \code{N\_RUNS} volte,
aggrega i risultati, e --- se la distanza scende sotto
\code{CONFIRM\_BELOW} --- aggiunge \code{CONFIRM\_EXTRA\_RUNS} run di
\emph{conferma anti-rumore} (le simulazioni sono stocastiche, non ci si
fida di un singolo ``colpo fortunato''). Alla fine:
\begin{itemize}
  \item \code{fitness = min\_distance} $=$ minima distanza osservata;
  \item \code{avg\_point} $=$ media dei punti Formula 1 sui run;
  \item \code{avg\_flight\_min} $=$ durata media di volo;
  \item \code{traj} $=$ \emph{media} delle traiettorie ricampionate.
\end{itemize}

\subsection{Il ciclo principale \texttt{generate}}
È il metodo che orchestra tutto. In pseudocodice ad alto livello:
\begin{lstlisting}[style=pseudo]
funzione generate(budget):
    total_budget = budget
    se esiste un checkpoint valido:
        ripristina lo stato (archivio, parent, sigma, rotta, ...)
    altrimenti:
        bootstrap della rotta se necessario
        prepara INIT_SEEDS semi; sigma = SIGMA_INIT; stagnation = 0
        salva checkpoint

    # --- fase di seeding (ripristinabile seme per seme) ---
    mentre restano semi and budget disponibile:
        valuta il seme; archivialo se e' un fallimento
        aggiorna il "miglior seme"; salva checkpoint
    parent = miglior seme

    # --- ciclo evolutivo (1+1) ---
    mentre budget_used < total_budget:
        generation += 1
        se c'era un test "pending" (crash): rifallo da capo
        altrimenti: child = _make_child(parent, sigma)
        salva checkpoint con pending = child     # PRIMA di valutare
        _evaluate(child)
        pending = None

        improved = (child.fitness < parent.fitness)
        aggiorna success_history con improved
        se improved: parent = child; stagnation = 0
        altrimenti:  stagnation += 1

        esito = _maybe_archive(child)
        se esito == ARCH_REDUNDANT: stagnation += 1   # extra

        sigma = _adapt_sigma(sigma, success_history)  # regola 1/5

        se stagnation >= STAGNATION_LIMIT:
            parent = _restart_individual()
            sigma = SIGMA_INIT; stagnation = 0; success_history = []
        salva checkpoint (fine generazione)

    ritorna _ranked_archive()   # ordinato per Formula 2
\end{lstlisting}

\subsection{L'operatore di mutazione \texttt{\_make\_child}}
\begin{lstlisting}[style=pseudo]
funzione _make_child(parent, sigma):
    nuovi_geni = [ _mutate_gene(g, sigma) per ogni g in parent.genes ]
    roll = numero casuale in [0,1)
    se roll < 0.10 and numero_geni < 3:          # MACRO: aggiungi
        extra = None
        se parent ha traiettoria and casuale < TRAP_PROB:
            extra = _trap_gene(parent, nuovi_geni)   # tenta il mirato
        se extra == None:                            # fallback (o prob non colpita
            extra = _non_overlapping_gene(nuovi_geni) #  o trap fallito): casuale
        se extra != None: aggiungi extra
    altrimenti se roll > 0.90 and numero_geni > 1: # MACRO: rimuovi
        rimuovi un gene a caso
    nuovi_geni = _repair_overlaps(nuovi_geni)
    ritorna Individual(nuovi_geni)
\end{lstlisting}
Due tipi di mutazione convivono: \textbf{micro} (rumore gaussiano su
tutti i geni, sempre) e \textbf{macro} (aggiunta/rimozione di un
ostacolo, $10\%+10\%$).

\subsection{Adattamento del passo \texttt{\_adapt\_sigma}}
Implementa la regola $1/5$ di Rechenberg su una finestra recente:
\begin{lstlisting}[style=pseudo]
funzione _adapt_sigma(sigma, storia):
    finestra = ultimi min(10, STAGNATION_LIMIT) esiti
    se la finestra non e' piena: ritorna sigma (invariato)
    tasso = successi / dimensione_finestra
    se tasso > 0.2: sigma = sigma * 1.22     # esplora di piu'
    altrimenti:     sigma = sigma / 1.22     # affina
    ritorna clamp(sigma, SIGMA_MIN, SIGMA_MAX)
\end{lstlisting}

\subsection{Restart: diversificare o intensificare}
Quando la ricerca stagna, \code{\_restart\_individual} sceglie una nuova
soluzione di partenza. La rotta è divisa in \emph{bucket} lungo l'asse
$y$ (larghi \code{Y\_BUCKET\_SIZE}).
\begin{lstlisting}[style=pseudo]
funzione _restart_individual():
    coperti = { bucket_y(ind) per ogni ind in archivio }
    vuoti   = bucket non coperti

    se archivio vuoto: ritorna un candidato casuale singolo
    se non ci sono bucket vuoti and casuale < INTENSIFY_PROB:
        ritorna _intensify_individual()      # INTENSIFICA
    # altrimenti DIVERSIFICA:
    genera RESTART_CANDIDATES candidati (preferendo bucket vuoti,
        restando sulla rotta)
    ritorna il candidato con IoU minima rispetto all'archivio
\end{lstlisting}
\begin{itemize}
  \item \textbf{Diversify} (default): riparte da una zona spaziale non
        ancora coperta, per aumentare la diversità dei fallimenti.
  \item \textbf{Intensify}: se tutte le zone sono coperte, riparte dal
        fallimento più severo in archivio e gli aggiunge un trap-gene,
        per spingerlo verso la collisione.
\end{itemize}

\subsection{La politica di archivio \texttt{\_maybe\_archive}}
È la logica di diversità, ``diversity-first'' nello spazio delle
traiettorie:
\begin{lstlisting}[style=pseudo]
funzione _maybe_archive(ind):
    se ind.min_distance >= SOFT_FAIL: ritorna ARCH_NONE  # non e' fallimento
    _dump_fail(ind)                    # salva SEMPRE su all_fails/

    dup = voce in archivio con traiettoria entro DTW_THRESHOLD (la piu' vicina)
    se dup != None:
        se ind e' piu' severo di dup:
            sostituisci dup con ind;  ritorna ARCH_IMPROVED
        altrimenti:                   ritorna ARCH_REDUNDANT
    se |archivio| < ARCHIVE_CAP:
        aggiungi ind;                 ritorna ARCH_NEW
    # archivio pieno: entra solo se batte il peggiore per Formula 2
    victim = voce con _ranking_score minimo
    se S(ind) <= S(victim):           ritorna ARCH_REDUNDANT
    sostituisci victim con ind;       ritorna ARCH_NEW
\end{lstlisting}
Tre casi, in ordine: stesso ``bacino'' di traiettoria (vince il più
severo); traiettoria nuova con slot libero; archivio pieno (competizione
per punteggio).

\subsection{Consegna: \texttt{\_delivery\_selection} e \texttt{write\_final\_selection}}
La consegna finale prende l'archivio ordinato per Formula 2 e lo filtra
\emph{greedy} per separazione DTW: si tiene un test solo se dista almeno
\code{DTW\_THRESHOLD} da tutti quelli già tenuti. Poi
\code{write\_final\_selection} copia i file da \code{all\_fails/} a
\code{consegna/} con nomi \code{rankNN\_*} (\code{rank01} = migliore).
L'operazione è \emph{idempotente} e non tocca mai \code{all\_fails/}.

\subsection{Crash-safety: il meccanismo di checkpoint}
Il sistema deve sopravvivere ai crash di Gazebo. La chiave è il campo
\code{pending}:
\begin{enumerate}
  \item \emph{prima} di valutare un candidato, i suoi geni vengono
        scritti nel checkpoint come \code{pending};
  \item il checkpoint è scritto in modo \textbf{atomico} (file temporaneo
        $+$ \code{os.replace}), così un crash a metà scrittura non
        corrompe l'ultimo checkpoint valido;
  \item al riavvio, se esiste un \code{pending}, quel candidato viene
        rieseguito \emph{da capo} (mai a metà);
  \item l'archivio e il budget non vengono mai resettati.
\end{enumerate}
Invariante di recupero: \emph{a ogni riavvio, tutti i fallimenti
completati prima del pending sono su disco, e la prossima azione valida
è ripetere quel pending dall'inizio}.

%------------------------------------------------------------
\section{\texttt{random\_generator.py} --- la baseline}
\label{sec:random}
%------------------------------------------------------------

\subsection{Ruolo}
Serve come \emph{gruppo di controllo} per l'esperimento: isola il
contributo della strategia evolutiva. Eredita \emph{tutto} da
\code{ESGenerator} (spazio legale, valutazione, archivio, checkpoint,
consegna) e ridefinisce solo due metodi.
\begin{lstlisting}[style=py]
class RandomGenerator(ESGenerator):
    def _sample_individual(self):
        n = random.randint(1, 3)          # 1..3 ostacoli
        genes = [_random_gene()]
        for _ in range(n - 1):
            extra = _non_overlapping_gene(genes)
            if extra is None: break
            genes.append(extra)
        return Individual(genes)

    def _make_child(self, parent, sigma):  # ignora parent e sigma
        return self._sample_individual()

    def _restart_individual(self):
        return self._sample_individual()
\end{lstlisting}
È l'applicazione dei pattern \textbf{Template Method} e \textbf{Strategy}:
\code{generate()} definisce il processo invariante, i metodi ridefiniti
sono i \emph{punti di variazione}. Poiché tutto il resto è identico, le
differenze nei risultati sono attribuibili \emph{solo} alla strategia di
ricerca.

%------------------------------------------------------------
\section{\texttt{cli.py} --- punto di ingresso}
\label{sec:cli}
%------------------------------------------------------------

\subsection{Ruolo e flusso}
È il \code{main} richiesto dalla competizione. Molto snello:
\begin{lstlisting}[style=pseudo]
GENERATOR = env("GENERATOR", "es")
se GENERATOR == "random": Generator = RandomGenerator
altrimenti:               Generator = ESGenerator

funzione main:
    configura i logger (debug su file, info su console)
    args = parse("generate <test> <budget>")
    generator = Generator(case_study_file = args.test)
    crea la cartella di output
    generator.generate(args.budget)          # la campagna
    generator.write_final_selection()        # copia in consegna/
    se eccezione: logga e esci con codice != 0
\end{lstlisting}
La selezione del generatore via variabile d'ambiente realizza il
\textbf{contratto stabile} richiesto: la riga di comando non cambia mai,
cambia solo \code{GENERATOR}.

%------------------------------------------------------------
\section{\texttt{plot\_official.py} --- visualizzazione}
\label{sec:plot}
%------------------------------------------------------------
Genera per ogni coppia \code{<nome>.yaml} $+$ \code{<nome>.ulg} un grafico
\code{<nome>\_plot.png} con: a sinistra $x,y,z,\text{yaw}$ nel tempo, a
destra la vista dall'alto con traiettoria e ostacoli ruotati. Legge la
traiettoria con \code{pyulog} (sull'host) o con Aerialist (fallback,
dentro il container). È volutamente \textbf{non critico}: se il plotting
fallisce, viene solo loggato e la ricerca non si ferma. Usa
\textbf{matplotlib} (grafici) e \textbf{pyyaml} (per leggere gli
ostacoli dal YAML).

%------------------------------------------------------------
\section{\texttt{start.sh} e \texttt{Dockerfile} --- deployment}
\label{sec:deploy}
%------------------------------------------------------------

\subsection{\texttt{start.sh}}
Wrapper Docker con \textbf{auto-resume}. Costruisce una cartella di
output specifica per la missione, pulisce container zombie, monta il
codice e la cartella di output, e avvia \emph{un solo} container per
tutto il budget. Se il container termina con codice diverso da zero
(crash), lo \textbf{riavvia} riprendendo dal checkpoint, fino a
\code{MAX\_RESTARTS} volte o fino a esaurimento budget. Un \code{trap} su
Ctrl+C ferma il ciclo di restart. Ripristina anche la proprietà
(ownership) dei file generati sull'host.

\subsection{\texttt{Dockerfile}}
Costruisce un'immagine self-contained estendendo
\code{skhatiri/aerialist}, installando le dipendenze di
\code{requirements.txt}, copiando il generatore e impostando l'agente di
default (\code{AGENT local}) e le risorse di PX4-Avoidance.

%============================================================
\chapter{Flusso di esecuzione completo}
%============================================================

\section{Diagramma di sequenza (una campagna intera)}
\begin{lstlisting}[style=ascii]
Utente   cli.py    ESGenerator   TestCase   Aerialist   PX4/Gazebo   FS
  |        |            |            |           |           |         |
  |--generate mission budget------->|           |           |         |
  |        |--init(case_study)----->|           |           |         |
  |        |            |--_load_route()------------------------------>| (.ulg)
  |        |            |--_set_step_scale()                 |         |
  |        |--generate(budget)----->|           |           |         |
  |        |            |--(no ckpt) seeding: valuta INIT_SEEDS        |
  |        |            |                                              |
  |        |     LOOP mentre budget disponibile:                      |
  |        |            |--_make_child(parent,sigma)                   |
  |        |            |--save_checkpoint(pending=child)------------->|
  |        |            |--create TestCase(obstacles)->|   |           |
  |        |            |--execute()------------------>|-->|           |
  |        |            |            |                 |   |--run mission->|
  |        |            |            |                 |   |<--trajectory--|
  |        |            |<--min_distance (get_distances)|              |
  |        |            |  (opz.) conferma se dist < CONFIRM_BELOW      |
  |        |            |--_maybe_archive (DTW dedup)                   |
  |        |            |  se fallimento: _dump_fail -> all_fails/----->|
  |        |            |--_adapt_sigma; restart/intensify se stagna    |
  |        |            |--save_checkpoint (fine generazione)---------->|
  |        |            |  [se crash: start.sh riavvia, pending da capo]|
  |        |            |                                              |
  |        |--write_final_selection()------------------------------->  |
  |        |            |--_delivery_selection (separazione DTW)        |
  |        |            |--copia rank01..rankNN in consegna/----------->|
  |<--log finale: "Delivery: N test in consegna/"-------              |
\end{lstlisting}

\section{Flusso decisionale del ciclo evolutivo}
\begin{lstlisting}[style=ascii]
             [nuova generazione]
                    |
          pending da crash? --si--> rifai lo stesso test da capo
                    | no
                    v
            _make_child(parent, sigma)
                    |
          salva checkpoint (pending = child)
                    |
                _evaluate(child)
                    |
        child.fitness < parent.fitness ?
            /                    \
          si                      no
           |                       |
   parent = child            stagnation += 1
   stagnation = 0                  |
           \                      /
            +--------+-----------+
                     v
              _maybe_archive(child)
                     |
        esito == ARCH_REDUNDANT ? --si--> stagnation += 1 (extra)
                     | no
                     v
             _adapt_sigma (regola 1/5)
                     |
        stagnation >= STAGNATION_LIMIT ?
            /                        \
          no                          si
           |                           |
           |               bucket y tutti coperti?
           |                   /              \
           |                 no               si, con prob INTENSIFY_PROB
           |                  |                     |
           |            DIVERSIFY               INTENSIFY
           |          (bucket vuoto,        (miglior fail + trap-gene)
           |           min IoU)                    |
           |                  \                    /
           |                   +--- reset sigma --+
            \                          /
             +--- salva checkpoint ---+
                        |
                  [prossima generazione]
\end{lstlisting}

%============================================================
\chapter{Dipendenze tra i moduli e librerie esterne}
%============================================================

\section{Dipendenze interne}
\begin{longtable}{L{3.6cm} L{10.4cm}}
\toprule
\textbf{Modulo} & \textbf{Importa da}\\
\midrule
\endhead
\code{cli.py} & \code{es\_generator} \emph{oppure} \code{random\_generator}.\\
\code{es\_generator.py} & \code{config}, \code{geometry}, \code{models}, \code{testcase}, (lazy) \code{plot\_official}.\\
\code{random\_generator.py} & \code{es\_generator}, \code{geometry}, \code{models}.\\
\code{geometry.py} & \code{config}, Aerialist (\code{Obstacle}, \code{Trajectory}).\\
\code{models.py} & \code{geometry} (\code{\_gene\_to\_obstacle}), \code{testcase}.\\
\code{testcase.py} & Aerialist (\code{AerialistTest}, agenti, \code{Trajectory}).\\
\bottomrule
\end{longtable}

\section{Librerie esterne e loro ruolo}
\begin{longtable}{L{3.4cm} L{10.6cm}}
\toprule
\textbf{Libreria} & \textbf{Ruolo nel progetto}\\
\midrule
\endhead
\textbf{Aerialist} & Banco di prova: carica la missione, crea gli ostacoli, esegue la simulazione PX4/Gazebo, estrae traiettoria e log. Fornisce anche gli oggetti \code{Obstacle} (geometria via Shapely).\\
\textbf{NumPy} & Aggregazione numerica: media delle traiettorie, matrici del DTW, ricampionamento.\\
\textbf{Matplotlib} & Generazione dei grafici (\code{plot\_official.py}).\\
\textbf{PyYAML} & Lettura degli ostacoli dai file YAML nel plotting.\\
\textbf{Shapely} & Operazioni geometriche (aree, intersezioni per IoU), usato tramite gli oggetti Aerialist.\\
\textbf{pyulog} & Lettura diretta dei log \code{.ulg} sull'host (fallback ad Aerialist nel container).\\
\textbf{python-decouple} & Lettura della variabile \code{AGENT} in \code{testcase.py} (dipendenza transitiva di Aerialist).\\
\textbf{Docker} & Deployment di riferimento e wrapper di riavvio su crash.\\
\bottomrule
\end{longtable}

%============================================================
\chapter{Riepilogo degli algoritmi principali}
%============================================================

\section{Obiettivo di ricerca}
\[
\min_{x\in\mathcal{X}} f(x), \qquad f(x)=d_{\min}(x).
\]
Scalare, non regolarizzato (indipendente da numero di ostacoli e tempo).
Ordinamento di severità: $f(x_1)<f(x_2)\Rightarrow x_1$ più critico.

\section{Tabella degli esiti di archivio}
\begin{longtable}{L{3.4cm} L{2cm} L{8.6cm}}
\toprule
\textbf{Costante} & \textbf{Valore} & \textbf{Significato}\\
\midrule
\endhead
\code{ARCH\_NONE} & 0 & Non è un fallimento (dist $\ge$ \code{SOFT\_FAIL}).\\
\code{ARCH\_NEW} & 1 & Traiettoria nuova: l'archivio è cresciuto.\\
\code{ARCH\_IMPROVED} & 2 & Stesso bacino, test migliore: rimpiazzata la voce.\\
\code{ARCH\_REDUNDANT} & 3 & Stesso bacino, non migliore: scartato (stagnazione extra).\\
\bottomrule
\end{longtable}

\section{Verdetto ufficiale (oracolo)}
\begin{longtable}{L{2.6cm} L{4.2cm} L{2cm} L{4.4cm}}
\toprule
\textbf{Verdetto} & \textbf{Condizione} & \textbf{Punti} & \textbf{Interpretazione}\\
\midrule
\endhead
Hard fail & $d_{\min}<0.25$ m & 5 & collisione o quasi-collisione\\
Soft fail & $0.25\le d_{\min}<1.0$ m & 2 & passaggio ravvicinato pericoloso\\
Soft fail & $1.0\le d_{\min}<1.5$ m & 1 & violazione marginale\\
No fail & $d_{\min}\ge 1.5$ m & 0 & volo sicuro / test rifiutato\\
\bottomrule
\end{longtable}

%============================================================
\chapter{Come modificare o estendere il progetto}
%============================================================

\section{Ricette pratiche}
\begin{description}
  \item[Aggiungere un nuovo generatore] Creare una sottoclasse di
    \code{ESGenerator} (come \code{RandomGenerator}) e ridefinire i punti
    di variazione (\code{\_make\_child}, eventualmente
    \code{\_restart\_individual}). Aggiungere un ramo in \code{cli.py}
    sulla variabile \code{GENERATOR}. Tutta l'infrastruttura (valutazione,
    archivio, checkpoint, consegna) si eredita gratis.
  \item[Cambiare la funzione di fitness] Modificare
    \code{\_evaluate} in \code{es\_generator.py} (dove si assegna
    \code{ind.fitness}). Attenzione: la fitness guida la ricerca, il
    \code{\_ranking\_score} guida solo la consegna; sono
    volutamente distinti.
  \item[Cambiare il criterio di diversità] Agire su \code{\_dtw\_mean} /
    \code{\_traj\_gap} (metrica) o su \code{DTW\_THRESHOLD} (soglia) in
    \code{config.py}.
  \item[Cambiare i vincoli sugli ostacoli] Modificare \code{MIN\_POS},
    \code{MAX\_POS}, \code{MIN\_SIZE}, \code{MAX\_SIZE} in
    \code{config.py}; le validità (\code{\_fit\_center},
    \code{\_boxes\_overlap}) si adeguano automaticamente.
  \item[Aggiungere un nuovo operatore di mutazione] Scriverlo in
    \code{geometry.py} (funzione pura, così è testabile senza
    simulatore) e richiamarlo da \code{\_make\_child}.
\end{description}

\section{Consigli di testabilità}
Grazie alla separazione \emph{Functional Core / Imperative Shell}, si può
testare quasi tutto senza avviare Gazebo: geometria, scoring,
serializzazione, DTW, logica di archivio e selezione sono funzioni
deterministiche. Impostando \code{SEED}, anche campionamento e mutazione
diventano ripetibili.

%============================================================
\chapter{Riassunto finale}
%============================================================

\section{I concetti da portare a casa}
\begin{enumerate}
  \item Il progetto è un \textbf{generatore di test evolutivo}: cerca
        configurazioni di ostacoli che minimizzano la distanza
        drone--ostacolo.
  \item Usa una \textbf{Strategia Evolutiva $(1+1)$}: un genitore, un
        figlio per generazione, selezione elitaria, passo adattivo
        (Rechenberg $1/5$).
  \item Campiona gli ostacoli \textbf{lungo la rotta} del drone e usa il
        \textbf{trap-gene} per bloccarne i varchi di fuga.
  \item Conserva solo fallimenti con \textbf{traiettorie diverse}
        (diversità DTW), separando \code{all\_fails/} (storia completa)
        da \code{consegna/} (sottoinsieme valutato).
  \item È \textbf{crash-safe} tramite checkpoint atomici e riesecuzione
        del candidato \code{pending}.
  \item La \textbf{fitness di ricerca} (distanza minima) è tenuta
        separata dal \textbf{ranking di consegna} (Formula 2), per non
        distorcere la ricerca con preferenze di semplicità.
\end{enumerate}

\section{Mappa mentale finale}
\begin{lstlisting}[style=ascii]
  PROBLEMA:  trovare ostacoli che fanno sbagliare il drone
      |
  STRATEGIA: (1+1)-ES che minimizza d_min
      |
  +--- init lungo la rotta (INIT_SEEDS semi)
  +--- mutazione: micro (gaussiana) + macro (aggiungi/rimuovi + trap-gene)
  +--- selezione: il figlio vince solo se piu' severo
  +--- passo adattivo: regola 1/5 di Rechenberg
  +--- stagnazione: restart (diversify) o intensify
      |
  ARCHIVIO: diversity-first per traiettoria (DTW), cap = ARCHIVE_CAP
      |
  CONSEGNA: sottoinsieme separato per DTW, ordinato per Formula 2
      |
  ROBUSTEZZA: checkpoint atomico + auto-resume Docker
\end{lstlisting}

\vfill
\begin{center}
\emph{Fine del manuale tecnico.}
\end{center}

\end{document}
